\documentclass[11pt,a4paper]{article}
\usepackage[margin=1in]{geometry}
\usepackage{mathptmx}
\usepackage{amsmath}
\usepackage{amssymb}
\usepackage{amsthm}
\usepackage{braket}
\usepackage{tikz}
\usepackage{quantikz}
\usepackage{enumitem}
\usepackage{hyperref}
\usepackage{xcolor}

\definecolor{circuitblue}{RGB}{40,80,160}
\hypersetup{colorlinks=true,linkcolor=circuitblue,urlcolor=circuitblue,citecolor=circuitblue}

\theoremstyle{definition}
\newtheorem{definition}{Definition}[section]
\newtheorem{example}{Example}[section]

\theoremstyle{remark}
\newtheorem*{intuition}{Intuition}
\newtheorem*{keypoint}{Key Point}

\newcounter{exercise}[section]
\newenvironment{exercise}{\refstepcounter{exercise}\par\medskip\noindent\textbf{Exercise \theexercise.}\rmfamily}{\medskip}

\pagestyle{plain}
\setlength{\parindent}{0pt}
\setlength{\parskip}{6pt}

\begin{document}

{\LARGE \textbf{Lab 1: Quantum Circuits}}\\[8pt]
{\large QML}\\[12pt]
\hrule

\vspace{16pt}

\section{introduction}

welcome to the first lab! in this session u will get ur hands dirty building quantum circuits from scratch using qiskit. the goal is to get comfortable with the basic workflow: create a circuit, add gates, simulate it, and look at the results.

by the end of this lab u should be able to:
\begin{itemize}[leftmargin=*]
  \item create and manipulate quantum circuits in qiskit
  \item build entangled states (bell states, GHZ states)
  \item implement quantum teleportation
  \item measure qubits in different bases
  \item use the statevector simulator and the qasm simulator
\end{itemize}

\begin{keypoint}
if u haven't installed qiskit yet, run \texttt{pip install qiskit} in ur terminal. we r using qiskit 1.x for this course. make sure u also have \texttt{qiskit-aer} installed for the simulators.
\end{keypoint}

\section{qiskit basics}

\subsection{the workflow}

every qiskit program follows the same basic pattern. here's wt u need to do:

\begin{enumerate}[leftmargin=*]
  \item import \texttt{QuantumCircuit} from \texttt{qiskit}
  \item create a circuit object --- u specify how many qubits and how many classical bits u want
  \item add gates to the circuit (hadamard, CNOT, etc.)
  \item choose a backend (simulator or real device)
  \item transpile and run
  \item grab the results
\end{enumerate}

the two simulators we care about r:
\begin{itemize}[leftmargin=*]
  \item \textbf{StatevectorSimulator} --- gives u the full quantum state. great for debugging bcs u can see exactly wt the state looks like. not realistic tho, bcs on a real quantum computer u can't peek at the state.
  \item \textbf{QasmSimulator} --- simulates actual measurements. u get a dictionary of counts, just like u would from a real device.
\end{itemize}

\subsection{single qubit gates}

before we start building circuits, let's review the gates we'll use. u should already know these from lecture but here's a quick refresher.

\begin{definition}
the \textbf{Hadamard gate} $H$ maps $\ket{0} \to \frac{1}{\sqrt{2}}(\ket{0}+\ket{1})$ and $\ket{1} \to \frac{1}{\sqrt{2}}(\ket{0}-\ket{1})$. in matrix form:
\[
H = \frac{1}{\sqrt{2}}\begin{pmatrix} 1 & 1 \\ 1 & -1 \end{pmatrix}
\]
\end{definition}

the Pauli gates r also important:
\[
X = \begin{pmatrix} 0 & 1 \\ 1 & 0 \end{pmatrix}, \quad
Y = \begin{pmatrix} 0 & -i \\ i & 0 \end{pmatrix}, \quad
Z = \begin{pmatrix} 1 & 0 \\ 0 & -1 \end{pmatrix}
\]

and for rotation we have $R_y(\theta)$ and $R_z(\theta)$ which will be super important later when we get to variational circuits.

\subsection{two qubit gates}

the big one is CNOT (also called CX). it flips the target qubit if the control qubit is $\ket{1}$.

here's wt a simple circuit with H and CNOT looks like --- this is the circuit that creates a bell state:

\begin{center}
\begin{quantikz}
\lstick{$q_0$} & \gate{H} & \ctrl{1} & \meter{} \\
\lstick{$q_1$} & \qw & \targ{} & \meter{}
\end{quantikz}
\end{center}

\begin{intuition}
the H gate puts $q_0$ into superposition, and then the CNOT entangles $q_0$ and $q_1$. the result is $\frac{1}{\sqrt{2}}(\ket{00}+\ket{11})$, which is one of the four bell states. when u measure, u'll get either $00$ or $11$ with equal probability --- never $01$ or $10$. that's entanglement for u.
\end{intuition}

\begin{exercise}
\textbf{ur first circuit.}
create a quantum circuit with 2 qubits and 2 classical bits. apply a hadamard gate to qubit 0, then a CNOT with control on qubit 0 and target on qubit 1. add measurements to both qubits.

steps:
\begin{enumerate}[leftmargin=*]
  \item import \texttt{QuantumCircuit} from qiskit
  \item create \texttt{QuantumCircuit(2, 2)}
  \item call \texttt{.h(0)} to add hadamard on qubit 0
  \item call \texttt{.cx(0, 1)} for the CNOT
  \item call \texttt{.measure([0,1], [0,1])} to measure both qubits
  \item draw the circuit using \texttt{.draw('mpl')}
\end{enumerate}

run this on the qasm simulator with 1024 shots. plot a histogram of the results. u should see roughly 50/50 split between \texttt{00} and \texttt{11}.
\end{exercise}

\section{bell states}

there r four bell states and they form a basis for the two-qubit space. every bell state is maximally entangled.

\begin{align}
\ket{\Phi^+} &= \frac{1}{\sqrt{2}}(\ket{00}+\ket{11}) \\
\ket{\Phi^-} &= \frac{1}{\sqrt{2}}(\ket{00}-\ket{11}) \\
\ket{\Psi^+} &= \frac{1}{\sqrt{2}}(\ket{01}+\ket{10}) \\
\ket{\Psi^-} &= \frac{1}{\sqrt{2}}(\ket{01}-\ket{10})
\end{align}

the trick to building all four is to start with the right input state and then apply H + CNOT.

circuit for $\ket{\Phi^-}$:

\begin{center}
\begin{quantikz}
\lstick{$q_0$} & \gate{X} & \gate{H} & \ctrl{1} & \qw \\
\lstick{$q_1$} & \qw & \qw & \targ{} & \qw
\end{quantikz}
\end{center}

circuit for $\ket{\Psi^+}$:

\begin{center}
\begin{quantikz}
\lstick{$q_0$} & \gate{H} & \ctrl{1} & \qw \\
\lstick{$q_1$} & \gate{X} & \targ{} & \qw
\end{quantikz}
\end{center}

\begin{exercise}
\textbf{all four bell states.}
build circuits for all four bell states. for each one:
\begin{enumerate}[leftmargin=*]
  \item construct the circuit (no measurements yet)
  \item use the statevector simulator to get the output state
  \item verify the state vector matches the expected bell state
  \item then add measurements and run on qasm simulator with 4096 shots
  \item record the measurement outcomes
\end{enumerate}

hint: $\ket{\Phi^+}$ is just H + CNOT on $\ket{00}$. for $\ket{\Phi^-}$, apply X to qubit 0 first. for $\ket{\Psi^+}$, apply X to qubit 1 first. for $\ket{\Psi^-}$, apply X to both qubits first.

write down the mapping: which input state $\to$ which bell state?
\end{exercise}

\section{GHZ states}

a GHZ state is the generalization of a bell state to $n$ qubits:
\[
\ket{\text{GHZ}_n} = \frac{1}{\sqrt{2}}(\ket{0}^{\otimes n} + \ket{1}^{\otimes n})
\]

for 3 qubits that's $\frac{1}{\sqrt{2}}(\ket{000}+\ket{111})$. the circuit is pretty straightforward:

\begin{center}
\begin{quantikz}
\lstick{$q_0$} & \gate{H} & \ctrl{1} & \qw & \qw \\
\lstick{$q_1$} & \qw & \targ{} & \ctrl{1} & \qw \\
\lstick{$q_2$} & \qw & \qw & \targ{} & \qw
\end{quantikz}
\end{center}

\begin{intuition}
the pattern is: H on the first qubit, then a chain of CNOTs. each CNOT spreads the entanglement to one more qubit. for $n$ qubits u need 1 H gate and $n-1$ CNOT gates.
\end{intuition}

\begin{exercise}
\textbf{building GHZ states.}
\begin{enumerate}[leftmargin=*]
  \item build a 3-qubit GHZ circuit. verify the statevector is $\frac{1}{\sqrt{2}}(\ket{000}+\ket{111})$.
  \item build a 5-qubit GHZ circuit. measure it and check u only get $\ket{00000}$ and $\ket{11111}$.
  \item write a function that takes $n$ as input and returns an $n$-qubit GHZ circuit. the function should:
    \begin{itemize}
      \item create a \texttt{QuantumCircuit(n)}
      \item apply H to qubit 0
      \item loop from $i=0$ to $n-2$ and apply CNOT with control $i$ and target $i+1$
      \item return the circuit
    \end{itemize}
  \item test ur function for $n = 2, 3, 4, 5, 8$. for each one, verify with statevector simulator.
\end{enumerate}
\end{exercise}

\section{measurement in different bases}

so far we've only measured in the computational basis ($Z$ basis). but sometimes u want to measure in other bases, like the $X$ basis or the $Y$ basis.

the trick: u can't directly measure in a different basis in qiskit. instead, u apply a rotation before measuring that transforms the basis u care about into the computational basis.

\begin{itemize}[leftmargin=*]
  \item \textbf{$X$ basis measurement}: apply H before measuring. this works bcs H maps $\ket{+} \to \ket{0}$ and $\ket{-} \to \ket{1}$.
  \item \textbf{$Y$ basis measurement}: apply $S^\dagger$ then H before measuring.
\end{itemize}

circuit for measuring in the $X$ basis:

\begin{center}
\begin{quantikz}
\lstick{$q_0$} & \qw & \gate{H} & \meter{}
\end{quantikz}
\end{center}

circuit for measuring in the $Y$ basis:

\begin{center}
\begin{quantikz}
\lstick{$q_0$} & \qw & \gate{S^\dagger} & \gate{H} & \meter{}
\end{quantikz}
\end{center}

\begin{exercise}
\textbf{basis measurements.}
\begin{enumerate}[leftmargin=*]
  \item prepare the state $\ket{+} = \frac{1}{\sqrt{2}}(\ket{0}+\ket{1})$ by applying H to $\ket{0}$.
  \item measure it in the $Z$ basis. wt do u get? (should be 50/50)
  \item measure it in the $X$ basis (apply H before measurement). wt do u get now? (should be 100\% one outcome)
  \item explain why this makes sense.
  \item now prepare $\ket{-} = \frac{1}{\sqrt{2}}(\ket{0}-\ket{1})$ (apply X then H to $\ket{0}$). measure in the $X$ basis. how does this differ from measuring $\ket{+}$?
\end{enumerate}
\end{exercise}

\section{quantum teleportation}

this is the big one for this lab. quantum teleportation lets u transfer a quantum state from one qubit to another using entanglement and classical communication.

the protocol:
\begin{enumerate}[leftmargin=*]
  \item alice and bob share a bell state $\ket{\Phi^+}$
  \item alice has a qubit in some unknown state $\ket{\psi} = \alpha\ket{0}+\beta\ket{1}$ that she wants to send to bob
  \item alice applies CNOT (her qubit controlling the bell state qubit) then H on her qubit
  \item alice measures her two qubits
  \item based on the measurement results, bob applies corrections (X and/or Z gates)
  \item bob's qubit is now in state $\ket{\psi}$
\end{enumerate}

here's the full circuit:

\begin{center}
\begin{quantikz}
\lstick{$\ket{\psi}$} & \qw & \qw & \ctrl{1} & \gate{H} & \meter{} & \qw & \qw \\
\lstick{$q_1$} & \gate{H} & \ctrl{1} & \targ{} & \qw & \meter{} & \qw & \qw \\
\lstick{$q_2$} & \qw & \targ{} & \qw & \qw & \qw & \gate{X} & \gate{Z}
\end{quantikz}
\end{center}

the first two gates (H on $q_1$, CNOT on $q_1 \to q_2$) create the bell pair. then the CNOT and H on the top two qubits r alice's operations. the X and Z on the bottom qubit r bob's corrections (controlled by alice's measurement results).

\begin{exercise}
\textbf{quantum teleportation.}
\begin{enumerate}[leftmargin=*]
  \item build the teleportation circuit step by step:
    \begin{itemize}
      \item create a 3-qubit circuit with 2 classical bits (for alice's measurements)
      \item prepare the state to teleport on qubit 0 --- use $R_y(\pi/4)$ to make smtg interesting
      \item save the statevector of qubit 0 (u'll compare later)
      \item create the bell pair between qubits 1 and 2
      \item apply CNOT(0,1) then H(0)
      \item measure qubits 0 and 1
      \item apply conditional X and Z gates on qubit 2
    \end{itemize}
  \item for the conditional gates, in qiskit u can use \texttt{.c\_if()} or the newer \texttt{if\_test} context manager. describe wt each classical bit controls.
  \item verify that after teleportation, qubit 2 ends up in the same state as qubit 0 started in. u can do this by:
    \begin{itemize}
      \item running many shots and comparing measurement statistics
      \item or using the statevector simulator (tho this is a bit tricky with mid-circuit measurements)
    \end{itemize}
  \item try teleporting different states: $\ket{0}$, $\ket{1}$, $\ket{+}$, $R_y(\pi/3)\ket{0}$.
\end{enumerate}
\end{exercise}

\section{putting it all together}

\begin{exercise}
\textbf{circuit challenge.}
build a circuit that:
\begin{enumerate}[leftmargin=*]
  \item creates a 4-qubit GHZ state
  \item measures qubits 0 and 1 in the $X$ basis
  \item measures qubits 2 and 3 in the $Z$ basis
  \item runs on the qasm simulator with 8192 shots
\end{enumerate}

analyze the correlations between the measurement outcomes. r the $X$-basis results correlated with the $Z$-basis results? explain why or why not.

the circuit should look smtg like this:

\begin{center}
\begin{quantikz}
\lstick{$q_0$} & \gate{H} & \ctrl{1} & \qw & \qw & \gate{H} & \meter{} \\
\lstick{$q_1$} & \qw & \targ{} & \ctrl{1} & \qw & \gate{H} & \meter{} \\
\lstick{$q_2$} & \qw & \qw & \targ{} & \ctrl{1} & \qw & \meter{} \\
\lstick{$q_3$} & \qw & \qw & \qw & \targ{} & \qw & \meter{}
\end{quantikz}
\end{center}
\end{exercise}

\begin{exercise}
\textbf{bonus: superdense coding.}
superdense coding is the reverse of teleportation --- u send 2 classical bits using 1 qubit (plus a shared bell pair).

the protocol:
\begin{itemize}[leftmargin=*]
  \item alice and bob share $\ket{\Phi^+}$
  \item alice encodes her 2-bit message by applying I, X, Z, or XZ to her qubit
  \item alice sends her qubit to bob
  \item bob applies CNOT then H and measures both qubits to recover the message
\end{itemize}

implement this for all four possible messages (00, 01, 10, 11). verify bob always recovers the correct message.
\end{exercise}

\section{submission}

submit a jupyter notebook with:
\begin{itemize}[leftmargin=*]
  \item all exercises completed with code and output
  \item circuit diagrams (use \texttt{.draw('mpl')})
  \item brief explanations of wt u observe
  \item the teleportation circuit working for at least 2 different input states
\end{itemize}

\end{document}
